\section{Truth Table Strategies}

\subsection{Simplification Methods}

The truth table method is tedious to say the least. Even when the number of
atomic sentences makes the overall number of rows manageable, a logician would
need to write dozens of T and F symbols to completely fill in a table. We can
introduce a number of shortcuts to make the process much more ergonomic.

\begin{observation}[Methods to Simplify Truth Tables]
  \label{obs:tt-simplified-methods}
  \begin{enumerate}
  \item You do not need to repeat truth values for the atomic letters. They are
    easy enough to read right off the left column.
  \item Since \ephi\eor\epsi is true just in case either \ephi or \epsi is true,
    we can skip one disjunct if the other is already marked T.
  \item Since \ephi\eand\epsi is false just in case either \ephi or \epsi is
    false, we can skip one conjunct if the other is already marked F.
  \item Since \ephi\eif\epsi is true in case either \epsi is true or \ephi is
    false, we can skip one component if the other is already so marked.
  \end{enumerate}
\end{observation}

This truth table with four atomic sentences shows how much savings we can get
from the simplified approach. Not only does it have 16 rows, but the complexity
of the sentence involves quite a few columns too. % (p 109)

\begin{center}
  \begin{tabular}{c c c c|c c c c c c c c c c c c c c c c}
    P&Q&R&S&$(((S$&\eand&$Q)$&\eand&\enot&\enot&$R)$&\eor         &\enot&$(\enot$&$(P$&\eor&\enot&$P)$&\eand&$Q))$ \\
    \hline
    T&T&T&T&      &     &    &     &     &     &    &           T &  T  &        &    &    &     &    &   F & \\
    T&T&T&F&      &     &    &     &     &     &    &           T &     &        &    &    &     &    &     & \\
    T&T&F&T&      &     &    &     &     &     &    &           T &     &        &    &    &     &    &     & \\
    T&T&F&F&      &     &    &     &     &     &    &           T &     &        &    &    &     &    &     & \\
    T&F&T&T&      &     &    &     &     &     &    &           T &  T  &        &    &    &     &    &   F & \\
    T&F&T&F&      &     &    &     &     &     &    &           T &     &        &    &    &     &    &     & \\
    T&F&F&T&      &     &    &     &     &     &    &           T &     &        &    &    &     &    &     & \\
    T&F&F&F&      &     &    &     &     &     &    &           T &     &        &    &    &     &    &     & \\
    F&T&T&T&      &     &    &  T  &     &     &    &\underline{T}&  F  &        &    &    &     &    &   T & \\
    F&T&T&F&      &     &    &  F  &     &     &    &\underline{F}&     &        &    &    &     &    &     & \\
    F&T&F&T&      &     &    &  F  &     &     &    &\underline{F}&     &        &    &    &     &    &     & \\
    F&T&F&F&      &     &    &  F  &     &     &    &\underline{F}&     &        &    &    &     &    &     & \\
    F&F&T&T&      &     &    &     &     &     &    &           T &  T  &        &    &    &     &    &   F & \\
    F&F&T&F&      &     &    &     &     &     &    &           T &     &        &    &    &     &    &     & \\
    F&F&F&T&      &     &    &     &     &     &    &           T &     &        &    &    &     &    &     & \\
    F&F&F&F&      &     &    &     &     &     &    &           T &     &        &    &    &     &    &     & \\
    \hline
     & & & &      &     &    &  4  &     &     &    & 3/5         &  2  &        &    &    &     &    &   1 & \\
  \end{tabular}
\end{center}

We can summarize the steps of the reasoning as follows.

\begin{enumerate}
\item The chunking of truth values that we use in constructing the table
  exhaustively is really paying off here. Since P and Q are the first two
  sentences in the table, they have long runs of rows where the truth values of
  both remains unchanged. So whatever the value of $\enot(P \eor P) \eand Q$ is
  it will be the same for each four-row block and when we figure it out we can
  just write it down every four lines. In another lucky break, we know that
  $P \eor P$ is always just whatever P is. Taking this all together
  $\enot(P \eor P) \eand Q$ is true just in case P is false and Q is true: so
  only for the four rows in the third block.
\item The same reasoning about the chunking means that we can just reverse
  these values in the columns for $\enot(\enot(P \eor P) \eand Q)$: making all
  the F values T, and vice-versa.
\item Since the main connective is \eor, we know that only one of the disjuncts
  needs to be true. The right disjunct is true on 12 rows, so we can just fill
  those all in for the main connective now. We don't need to know the value of
  the left disjunct in these cases.
\item We need to compute the truth value of the left disjunct for the remaining
  4 rows. Noting the double negation $\enot\enot R$, we can see that the entire
  expression $(S \eand Q)\eand\enot\enot R$ is true just in case all of S, Q and
  R are. This is just one of the remaining rows. The rest are false.
\item Now we have all we need to fill in the last 4 rows for the main
  connective. These four truth values are \underline{underlined} in the truth
  table.
\end{enumerate}

% more examples (pp 132-3)

\subsection{Truth-Theoretic Concepts}

Now that we have a handle on the truth-functional semantics of TFL, we are
better positioned to put meat on the bones of some informal concepts introduced
a few weeks ago: necessary truth/falsehood, contingent truth/falsehood,
equivalence and joint possibility.

We can use truth tables to define a technical notions that approximate the
concepts of \emph{necessary truth} and \emph{necessary falsity}. First, we can
think of these concepts in terms of the forms of arguments. Let's start by
firming up our definition of \emph{argument} from week 1.

\begin{definition}[Argument]
  \label{def:tfl-argument}
  An \emph{argument} is a sequence of sentences
  $\ephi_1\dots\ephi_n \therefore \epsi$. It comprises some \emph{premises},
  ($\ephi_1\dots\ephi_n$) and a \emph{conclusion} (\epsi), separated by an
  inference marker (\therefore).
\end{definition}

Beware that not all arguments are \emph{valid}. We can make bad arguments and
checking to see if an argument is good is a big part of the point of our formal
apparatus. We just want to define the kinds of things that we will check for
validity. Indeed, eventually we'll just drop the reference to arguments from our
definition altogether.

\begin{definition}[Tautologically Valid]
  \label{def:tfl-tautologically-valid}
  A TFL argument $\ephi_1,\dots,\ephi_n \therefore \epsi$ is
  \emph{tautologically valid} (i.e. \emph{logically valid}) iff there is no
  valuation where all of $\ephi_1,\dots,\ephi_n$ are true but $\epsi$ is false.
\end{definition}

We can characterize the relationship between the premises and conclusion of
a tautologically valid argument as one of \emph{entailment}.

\begin{definition}[Tautologically Entail]
  \label{def:tfl-tautologically-entail}
  A collection of TFL sentences $\ephi_1,\dots,\ephi_n$ \emph{tautologically
  entails} (i.e. \emph{logically entails}) a TFL sentence \epsi iff
  $\ephi_1,\dots,\ephi_n \therefore \epsi$ is tautologically valid. When this
  is the case, we will use the symbol \eentails instead of \therefore.

  \begin{argument}
    \item[ ] $\ephi_1,\dots,\ephi_n \eentails \epsi$
  \end{argument}

  Otherwise, we will say that $\ephi_1,\dots,\ephi_n$ \emph{does not entail}
  \epsi and use the symbol \enotentails.

    \begin{argument}
    \item[ ] $\ephi_1,\dots,\ephi_n \enotentails \epsi$
    \end{argument}
\end{definition}

We can recast the notion of tautologocal validity to apply not just to
arguments, but to individual sentences as well. Naturally, the definition is
extremely similar.

\begin{definition}[Tautology]
  \label{def:tfl-tautology}
  \ephi is a \emph{tautology} in TFL iff it is true on every valuation. When
  a \ephi is a tautology, we will signify this by using \eentails with no
  premises.

  \begin{argument}
    \item[ ] \eentails\ephi
  \end{argument}

  Otherwise \ephi is not a tautology, which we can signify in the metalogic with
  \enotentails.
  
  \begin{argument}
    \item[ ] \enotentails\ephi
  \end{argument}
  
\end{definition}

Some well-known examples of a tautologies have come up already:

\begin{argument}
\item[\textbf{LEM:}] A\eor\enot A
\item[\textbf{LNC:}] \enot(A\eand\enot A)
\item[\textbf{DN:}]  \enot\enot A\eiff A  
\end{argument}

We can confirm that each is invariably true using simplified truth tables

\begin{table}[h]

  \threeitems{
  \begin{center}
    \begin{tabular}{c|c c c}
      A      & A     & \eor           & \enot A \\
      \hline
      T      &       & \textbf{T}     &      F \\
      F      &       & \textbf{T}     &      T    
    \end{tabular}
  \end{center}
  }{
  \begin{center}
    \begin{tabular}{c|c c c c}
      $A$& \enot      & $(A$ & \eand  & $\enot A)$ \\
      \hline
      T      & \textbf{T} &          & F      &     F \\
      F      & \textbf{T} &          & F      &     T
    \end{tabular}
  \end{center}
  }{
  \begin{center}
    \begin{tabular}{c|c c c c}
      A & \enot & \enot A & \eiff      & A  \\
      \hline
      T      &   T   &      F     & \textbf{T} & \\
      F      &   F   &      T     & \textbf{T} &
    \end{tabular}
  \end{center}
  }
  \label{tab:tt-neg-equivalences}
  \caption{Simplified truth tables demonstrating \textbf{LEM}, \textbf{LNC}, and \textbf{DN}}
\end{table}

Conversely, a sentence that is false on every valuation is a
\emph{contradiction}.

\begin{definition}[Contradiction]
  \label{def:tfl-contradiction}
  The TFL sentence \ephi is a \emph{contradiction} iff it is false on every
  valuation.
\end{definition}

For example, $A\eand\enot A$ is a contradiction whereas
$A\eor\enot A$ is a tautology.

\begin{table}[h]
  \begin{center}
    \begin{tabular}{c||c|c}
      $A$ & $A \eand \enot A$ & $A \eor \enot A$ \\
      \hline
      T & \textbf{F} & \textbf{T} \\
      F & \textbf{F} & \textbf{T}
    \end{tabular}
  \end{center}
  \label{tab:truth-table-example}
  \caption{A simplified truth table with a contradition and a tautology}
\end{table}

We will not add a symbol like \eentails that we can use for contradictions,
but we can leverage the fact that tautologies and contradictions are complements
to symbolize it in the metalogic anyway. If \ephi is a contradiction then
\eentails\enot\ephi and if \eentails\ephi then \enot\ephi is a contradiction.

A sentence that is neither a tautology nor a contradiction is called
\emph{contingent}.

\begin{definition}[Contingent Sentence]
  \label{def:tfl-contingent}
  A TFL sentence \ephi is a \emph{contingent} iff neither \eentails\ephi nor
  \eentails\enot\ephi.
\end{definition}

Similarly, we can test whether two sentences are equivalent by checking to
see if they have the same truth tables.

\begin{definition}[Equivalence]
  \label{def:tfl-equivalence}
  Two TFL sentences \ephi and \epsi are \emph{equivalent} in TFL iff \ephi and
  \epsi have the same truth value as each other in every valuation.
\end{definition}

So we can easily verify that $\enot(A\eand B)$ is equivalent to $\enot A \eor
\enot B$ since they have the same truth value in every row.

\begin{table}[h]
  \begin{center}
    \begin{tabular}{c c|c c c c|c c c c c}
      $A$&$B$ & \enot&$(A$&\eand&$B)$ & \enot&$A$&\eor&\enot&$B$ \\
      \hline
      T & T  &  \bf{F} & T & T & T  &  F & T & \bf{F} & F & T \\
      T & F  &  \bf{T} & T & F & F  &  F & T & \bf{T} & T & F \\
      F & T  &  \bf{T} & F & F & T  &  T & F & \bf{T} & F & T \\
      F & F  &  \bf{T} & F & F & F  &  T & F & \bf{T} & T & F
    \end{tabular}
  \end{center}
  \label{tab:tt-demorgans-nand}
  \caption{A complete truth table demonstrating the equivalence of
    $\enot(A\eand B)$ and $\enot A\eor\enot B$.}
\end{table}

Similarly, we can confirm that $\enot(A\eor B)$ is equivalent to
$\enot A \eand \enot B$.

\begin{table}[h]
  \begin{center}
    \begin{tabular}{c c|c c c c c|c c c}
      $A$&$B$&\enot  &$(A$&\eor&$B)$ & $\enot A$&$\eand$&$\enot B$ \\
      \hline
      T      & T     &\bf{F} &        & T  &         &  F          &\bf{F} & F \\
      T      & F     &\bf{F} &        & T  &         &  F          &\bf{F} & T \\
      F      & T     &\bf{F} &        & T  &         &  T          &\bf{F} & F \\
      F      & F     &\bf{T} &        & F  &         &  T          &\bf{T} & T
    \end{tabular}
  \end{center}
  \label{tab:tt-demorgans-nor}
  \caption{A simplified truth table demonstrating the equivalence of
    $\enot(A\eor B)$ and $\enot A\and\enot B$.}
\end{table}

%% NB: mention DeMorgans here?

Being able to show equivalence is a powerful capability. We've seen how we can
make use of it in the metalogic by abstracting away to \emph{forms} in our
rule schemas. We can also lean on the same inter-substitutability as we try to
assess the tautological status of various complex sentences.

\begin{observation}[Substitution instances and tautologies]
  If \eentails\ephi and \epsi is a substitution-instance of \ephi, then
  \eentails\epsi. But the same is not true of \enotentails\ephi.
\end{observation}

We can just as easily use truth table methods to evaluate the logical properties
of sets of sentences, including arguments.

\begin{definition}[Joint Satisfiability]
  \label{def:tfl-jointly-satisfiable}
  The sentences $\ephi_1\dots\ephi_n$ are \emph{jointly satisfiable} in TFL iff
  some valuation makes them all true. Otherwise, they are \emph{jointly
  unsatisfiable}.
\end{definition}

We can clean up our definition of \emph{entailment} now, too, dropping the
references to arguments and the inference marker (\therefore), which really
are not important to the concept as now understood.

\begin{definition}[Entailment]
  \label{def:tfl-entailment}
  The TFL sentences $\ephi_1\dots\ephi_n$ \emph{entail} the TFL sentence \epsi
  ($\ephi_1\dots\ephi_n\eentails\epsi$) iff no valuation makes
  $\ephi_1\dots\ephi_n$ all true but \epsi false.
\end{definition}

Now that we've dropped the references to arguments, we can look for symmetry in
entailment: \emph{equivalence}.

\begin{definition}[Mutual Entailment]
  \label{def:tfl-mutual-entailment}
  The TFL sentences \ephi and \epsi \emph{mutually entail} each other
  (\ephi\emutual\epsi) iff both \ephi\eentails\epsi and \epsi\eentails\ephi.
  \\
  When \ephi\emutual\epsi, we say that \ephi and \epsi are \emph{tautologically
  equivalent}.
\end{definition}

And as we've mentioned already, equivalence plays a fundamental role in formal
logic.

We've postponed a lot of discussion about the material conditional (\eif), the
limits of its interpretation, and its relationship to similar metalogical
notions like entailment. This is as good a time as any to wade into that
quagmire.

Let's start by (re-)defining the major components of a conditional.

\begin{definition}[Antecedent]
  If the material conditional \ephi\eif\epsi is a sentence of TFL, then \ephi
  is its \emph{antecedent}.
\end{definition}

\begin{definition}[Consequent]
  If the material conditional \ephi\eif\epsi is a sentence of TFL, then \epsi
  is its \emph{consequent}.
\end{definition}

We can also define some some basic structural relationships between material
conditionals, for example:

\begin{definition}[Converse]
  If the material conditional \ephi\eif\epsi is a sentence of TFL, then
  \epsi\eif\ephi is its \emph{converse}.
\end{definition}

\begin{definition}[Contrapositive]
  If the material conditional \ephi\eif\epsi is a sentence of TFL, then
  \enot\epsi\eif\enot\ephi is its \emph{contrapositive}.
\end{definition}

Likewise, we have some patterns of valid inference steps that amount to
hard-and-fast rules that we can apply without exception.

\begin{observation}[Modus Ponens (MP)]
  \label{obs:tfl-modus-ponens}
  If \ephi and \epsi are TFL sentences, then
  $ \epsi,\ephi\eif\epsi\eentails\epsi$ and this can be used as an inference
  step. This mode of inference is called \emph{modus ponens}.
\end{observation}

\begin{table}[h]
  \begin{center}
    \begin{tabular}{c c||c|c||c}
       A& B &  A\eif B &  A &  B \\
      \hline
      T & T  & T & T & T \\
      T & F  & F & T & F \\
      F & T  & T & F & T \\
      F & F  & T & F & T
    \end{tabular}
  \end{center}
  \label{tab:tt-modus-ponens}
  \caption{A truth table showing Modus Ponens is valid}
\end{table}

\begin{observation}[Modus Tollens (MT)]
  If \ephi and \epsi are TFL sentences, then
  $\enot\epsi,\ephi\eif\epsi\eentails\enot\ephi$ and this can be used as an
  inference step in TFL. This mode of inference is called \emph{modus tollens}.
\end{observation}

\begin{table}[h]
  \begin{center}
    \begin{tabular}{c c||c|c||c}
       A& B &  A\eif B & \enot B & \enot A \\
      \hline
      T    & T    &      T         &    F       &    F \\
      T    & F    &      F         &    T       &    F \\
      F    & T    &      T         &    F       &    T \\
      F    & F    &      T         &    T       &    T \\
    \end{tabular}
  \end{center}
  \label{tab:tt-modus-tollens}
  \caption{A truth table showing Modus Tollens is valid}
\end{table}

\begin{observation}[Falsehood Condition (FC)]
  If \epsi and \epsi are TFL sentences where \ephi is true and \epsi is false,
  then $\ephi\eif\epsi$ must be false.
\end{observation}

\begin{table}[h]
  \begin{center}
    \begin{tabular}{c c||c|c||c}
       A& B &  A\eif B \\
      \hline
      T    & F    &      F
    \end{tabular}
  \end{center}
  \label{tab:tt-falsehood-condition}
  \caption{A truth table showing the Falsehood Condition is valid}
\end{table}

\begin{observation}[Non-Reversability (NR)]
  If \ephi and \epsi are TFL sentences, then we cannot infer
   \ephi\eif\epsi from \ephi\eif\epsi . This is known as the
  \emph{non-reversability} of the material conditional.
\end{observation}

\begin{table}[h]
  \begin{center}
    \begin{tabular}{c c||c|c||c}
       A& B &  A\eif B &  B\eif A \\
      \hline
      T    & T    &        T       &        T  \\
      T    & F    &        F       &        T  \\
      F    & T    &\textbf{T}      &\textbf{F} \\
      F    & F    &        T       &        T
    \end{tabular}
  \end{center}
  \label{tab:tt-non-reversability}
  \caption{A truth table showing Non-Reversability is valid}
\end{table}

We've been mixing logic and metalogic in this lecture so haphazardly that it's
worth pausing here to make an important distinctions.

\begin{observation}[Entailment and material implication are truth-functionally equivalent]
   \ephi\eif\epsi iff  \ephi\eentails\epsi
\end{observation}

\begin{table}[h]
  \begin{center}
    \begin{tabular}{c c||c||c}
       A& B &  A\eif B &  A\eentails B \\
      \hline
      T    & T    &       T        &          T \\
      T    & F    &       F        &          F \\
      F    & T    &       T        &          T \\
      F    & F    &       T        &          T
    \end{tabular}
  \end{center}
  \label{tab:tt-non-reversability}
  \caption{A truth table showing Non-Reversability is valid}
\end{table}

Nevertheless, \eif is a connective of truth-functional logic and \eentails is
a symbol of the metalogic that we've built by augmenting English. We should
refrain from using one when we mean the other.

\subsection{Worries about the Material Conditional}

Now we are equipped with the vocabulary necessary to discuss the controversy
about material conditionals.

One way to think of these inference rules if pre-theoretically. In other words,
the reason that the entailments for \textbf{MP}, \textbf{MT}, \textbf{FC} and
\textbf{NR} come out the way they do is that we designed the truth table for
the material conditional (\eif) so that they would. In this light, we can view
them as \emph{assumptions} about how implication should work, given whatever
definition we decide to impose on the material conditional. They are working
as intended.

But one might still wonder if these assumptions are correct. Should we abandon
one or more of these principles in order to support a more ``natural'' type of
implication?

Most of the heartburn about the material conditional comes from the fact that
it is taken to be true whenever its antecedant is false. This has the effect
of making many conditionals \emph{vacuously} true: true \textbf{only} by virtue
of that fact.

\begin{argument}
  \item[ ] If Lions won the last Super Bowl, then pigs can fly.
\end{argument}

We know the truth of both antecedant (they did not) and consequent (they can
not), so we can calculate the truth of the conditional (it's true). But this
does not ``feel'' right. One might even argue that there is no situation in
which the consequent is true. Nevertheless ...

Let's try to firm up our discomfort with some help from
\cite[][pp. 166-167]{Smi20}. Recall that our definition of material implication
A \eif B is equivalent to \enot A\eor B. So when we allow the
inference in (D) we are also allowing the inferences in (M).

\threeitems{
  \begin{argument}
  \item[ ]\textbf{D} 
  \item[ ] \enot A
  \item[\therefore] \enot A\eor B
  \end{argument}
}{
  \begin{argument}
  \item[ ] \textbf{M1}
  \item[ ] \enot A
  \item[\therefore] A\eif B
  \end{argument}
}{
  \begin{argument}
  \item[ ] \textbf{M2}
  \item[ ] not A
  \item[\therefore] if A then B
  \end{argument}
}

But while (D) seems trivially correct, (M1) is not immediately obvious. And (M2)
comes across as just plain wrong.

But the Lions example is not so bad. The example is at least \emph{indicative},
meaning that its truth or falsity is fully settled and not hypothetical or
contingent on future events or the attitudes of human judges. We can live
with the idea that the conditional sentence is true, if only because its
implication will never come to pass. It is indeed \emph{vacuous}.

Suffice it to say, the material conditional just does not imply any sort of
\emph{causal connection} between its antecedent and its consequent. That's just
not part of how truth-functional logic works. Still, the material conditional
is a pretty good appoximation for what is actually said in an indicative
\emph{if...then} sentence of ordinary English, even if it leaves out most of
the pragmatic nuance. Its shortfalls are problably no worse than those of \eand
and \eor: ordinary English just behaves differently than TFL and that's OK.

Some evidence from ordinary English supports this acquiescence.

\begin{observation}[English conditionals are sometimes like A \eif B]
  \begin{outline}[enumerate]
    \1 English tautologies are typically equivalent to their
    contrapositives.
    \2 \emph{If I won the Powerball, then I am rich.}
    \2 \emph{If I am not rich, then I did not win the Powerball.}
    \1 The material conditional  A\eif B is just like \emph{if A then B}
    in some simple cases, and even correctly invalidates two common logical
    fallacies:
    \2 affirming the consequent: A\eif B,B \therefore A; and
    \2 denying the antecedent: A\eif B,\enot B \therefore \enot B.
    \1 Often \emph{if A then B} and \emph{A only if B} are equivalent, such as
    these examples (from \cite[][p. 155]{Smi20}):
    \2 \emph{Einstein's theory is right only if space-time is curved.}
    \2 \emph{If space-time isn't curved, then Einstein's theory isn't right.}
    \2 \emph{If Einstein's theory is right, then space-time is curved.}
  \end{outline}
\end{observation}

To see how indicative conditionals still have problems around the edges, it
helps to think about how we negate them. Consider the following symbolization
key.

\begin{argument}
\item[E:] God exists. 
\item[P:] God answers malevolent prayers.
\end{argument}

Putting these sentences together into a conditional ($E\eif P$) reads something
like (G).

\begin{argument}
\item[G:] If God exists, then God answers malevolent prayers.
\end{argument}

We can think of two ways of negating this sentence, but both are unsatisfying.

\begin{center}
  \begin{tabular}{|c|c|}
    \hline
    $\enot(E\eif P)$ & \makecell{It is not the case that if God exists\\God answers malevolent prayers.} \\
    \hline
    $E\eif\enot P$ & \makecell{If God exists, then God does not\\answer malevolent prayers.} \\
    \hline
  \end{tabular}
\end{center}

In the first example ($\enot(E\eif P)$), we would have to accept the truth of
$E$ (i.e. God exists), otherwise $G$ could not be false and so $\enot G$ could
not be true. In the second example ($E\eif\enot P$), the atheist is faced with
a different conundrum: when $E$ is false (i.e. when God does not exist), then
we can conclude \textbf{either} $P$ or $\enot P$, since $E\eif\enot P$ and
$E\eif P$ would then \textbf{both} be true.

Other the other end of the spectrum, we have \emph{counterfactuals} (i.e.
\emph{subjunctive} conditionals). David Lewis, perhaps the best known
philosopher of counterfactuals, suggests the following pair of examples
to illustrate the difference.\footnote{from \cite[][p 162]{Smi20})}

\begin{argument}
\item[I:] If Oswald didn't shoot Kennedy, then someone else did.
\item[S:] If Oswald hadn't shot Kennedy, then someone else would have.
\end{argument}

(I) basically just says that someone shot Kennedy. Of course, this seems more
or less right. In the case of (S), however, the sentence asks us to contemplate
a state of affairs that does not actually exist: the one where Oswald did not
shoot Kennedy. In \textbf{that} state of affairs, Kennedy may or may not have
been shot at all, which makes (S) less than a slam-dunk claim. This seems like
something we could debate, not something we should just take to be
\emph{vacuously true}, even if Oswald \textbf{did} shoot Kennedy.

Although the subjunctive/indicative division lines up well with the English
grammar conventions in the Oswald example, the alignment is not always
consistent (and mode is not always expressed using verbs). A better approach
is to think of the states-of-affairs as \emph{possible worlds} where we can
check what the truth-value of our counter-facutal sentence \textbf{would be} if
if it was evaluated in a world where the antecedant was true. This approach,
called \emph{intuitionistic logic}, is a type of \emph{modal logic} and
dissatisfaction with material implication was a major motive for its
formulation. But we will leave a survey of modal logic for Logic 2 and a
complete treatment for Modal Logic. For now we just note that possible worlds
semantics is not truth functional.

\begin{observation}
  The conditional expressed with a possible world semantics
  ($\ephi\Rightarrow\epsi$)is not truth-functional.
\end{observation}

To see this, we merely need to note that $\ephi\Rightarrow\epsi$ is true just
in case \epsi is true at the possible world where the sentence is evalauated.
Nothing about whether or not \ephi and \epsi are \emph{actually true} can help
us to figure this out.

Nevertheless, like democracy, the material conditional is the worst option
available, except for all the others. We just need to be careful to avoid
using it for the treatment of counterfactuals.

As we've already seen, the symbolization of English sentences into TFL leaves
much of the linguistic nuance out of the resulting TFL sentence. We can
summarize some of the important shortcomings of TFL symbolization as follows:

\begin{observation}[Shortcomings of TFL Symbolization]
  \label{obs:tfl-shortcomings}
  \begin{enumerate}
  \item TFL does not support the types of inferences depicted in Aristotles
    syllogisms (quantification).
  \item TFL insists on \textbf{LEM} (every TFL is either true
    or false), but some adjectives do not quite work like that.
  \item TFL comes with a number of ``paradoxes'' that make the material
    conditional unsatisfying for symbolizing English sentences: one the one
    hand, negating the whole conditional often gives the wrong conclusions
    about its antecedent sentence; on the other hand, negating the consequent
    sentence opens up the possibility of contradiction when the antecedant
    sentence is false.
  \end{enumerate}
\end{observation}

With the simiplified methods in \ref{obs:tt-simplified-methods}, we can much
more efficiently evaluate the validity of TFL arguments. Like the validity of
individual sentences, testing the validity of arguments depends on identifying
the load-bearing sentences in the structure of the argument. Tersely, we look
for \emph{bad lines}.

\begin{observation}[Bad Lines in a Truth Table for Entailment]
  \label{obs:tt-bad-lines}
  \begin{enumerate}
  \item When the conclusion is true, the line is not bad.
  \item When any of the premises is false, the line is not bad.
  \item Truth tables with no bad lines depict valid arguments.
  \end{enumerate}
\end{observation}

%TODO: Examples?? (p139ish?)

Last week, we saw how we can save time by only filling in the parts of our
truth tables that matter to the evaluation of the overall formula (i.e.
to its \emph{main operator}). But these are still \emph{complete} truth tables,
it's just that many of the cells are silent because they do not matter or are
obvious.

We can simplify our approach even more if we take a step back and consider what
we are trying to do in building truth tables in the first place. They are not
interesting in their own right. But they do give us a clear method for being
thorough in our exploration of the various logical possibilities, gven the rules
of the system.

One more step back. Why do we want to explore all of the logical possibilities
anyway? We want to be thorough in this way because we have defined a number of
our logical notions (e.g. \emph{tautology}, \emph{validity}, \emph{entailment})
in terms of these possibilities. When we explore these properties, we want to
show that the sentences being evaluated are true (or false) in all of the
possibilities. In our parlance from last week, we are looking for \emph{bad
lines} in the table.

But what about showing that a sentence \emph{does not} have these properties?
Surely it will suffice to just find one bad line, a \emph{counter-example} so
to speak. And indeed it will. Thus, we will refer to a such a bad line as a
\emph{contervaluation} We can describe the landscape with this table modified
from \cite{Mag25}[Chapter 15].

\begin{center}
  \begin{tabular}{|c|c|c|c|}
    \hline
                  & \bf{yes} & \bf{no} & should be (always)     \\
    \hline
    tautology     &  all     &  one    & true                   \\
    contradiction &  all     &  one    & false                  \\
    \hline
    equivalent    &  all     &  one    & same                   \\
    unsatisfiable &  all     &  one    & false                  \\
    \hline
    valid         &  all     &  one    & true conclusion or false premise  \\
    entailment    &  all     &  one    & true conclusion or false premise  \\
    \hline
  \end{tabular}
\end{center}

Recall from \ref{def:tfl-tautology} that a TFL sentence is a tautology just in
case it is true in all scenarios. Showing that a TFL sentence is not a
tautology amounts to finding a scenario where it is not true. Since we know
we are trying to find a bad (i.e. false) line, we work backward through the
columns with this assumption in mind.

\begin{center}
  \begin{tabular}{c c|c c c c c c c}
    A & B & (A & \eif   & B) & \eif   & (B & \eif & A)  \\
    \hline
    F & T &    & T      &    & \bf{F} &  T & F    & F   \\
    4 & 4 &    & 2      &    &   1    &  3 & 2    & 3   \\
  \end{tabular}
\end{center}

The TFL sentence $(A \eif B) \eif (B \eif A)$ is false when $A$ is false and
$B$ is true. It is not a tautology.

Similarly, we can show that a TFL sentence is not a contradiction by showing
that it is sometimes true. So this time we start by assuming a true line.

\begin{center}
  \begin{tabular}{c c|c c c c c c c}
    A & B & (A & \eif   & B) & \eif   & (B & \eif & A)  \\
    \hline
    T & F &  T & F      & F  & \bf{T} &    &      &     \\
    4 & 4 &  3 & 2      & 3  &   1    &    &      &     \\
  \end{tabular}
\end{center}

This time we had to decide what to do next, since there are two ways to make
the main operator true: let its antecedant be false or let its consequent be
true. We chose the former, because that choice (a false material implication)
forces our truth value assignment to the atomic letters: $A$ must be true and
$B$ must be false. We have found a counter-example. The TFL sentence
$(A\eif B)\eif(B\eif A)$ is also not a contradiction.

We can extend this practice to truth tables with multiple TFL sentences in
order to find counter-examples for equivalence, satisfiability, validity and
entailment. For example, the following partial table demonstrates that $A \eand
B$ and $A \eand \enot B$ are not equivalent.

\begin{center}
  \begin{tabular}{c c||c c c|c c c c}
    A & B    &    A & \eand  & B    &    A & \eand & \enot & B \\
    \hline
    T & T    &      & \bf{T} &      &      & \bf{F} &  F    &   \\
    2 & 2    &      &   1    &      &      &  4     &  3    &   \\
  \end{tabular}
\end{center}

When we set the values of $A$ and $B$ to make $A \eand B$ true, the truth value
of $A \eand \enot B$ correspondingly becomes false. The two TFL sentences have
different truth values in some scenarios. They cannot be equivalent. A similar
process can be applied to the remaining logical properties.

% TODO: satisfiability, validity and entailment tables ??

